\section{Vai trò của RAG trong hệ thống tư vấn tuyển sinh}

Trong hệ thống chatbot tư vấn tuyển sinh, thông tin cần cung cấp cho người dùng thường có phạm vi rộng, thay đổi theo từng năm và phụ thuộc vào nhiều loại tài liệu khác nhau như đề án tuyển sinh, thông báo tuyển sinh, quy chế xét tuyển, danh mục ngành đào tạo, phương thức xét tuyển, học phí, học bổng, lịch tuyển sinh và các câu hỏi thường gặp.

Nếu chỉ sử dụng các câu trả lời tĩnh được khai báo sẵn trong Rasa, hệ thống sẽ gặp khó khăn khi cần trả lời những câu hỏi có nội dung chi tiết, thay đổi thường xuyên hoặc yêu cầu trích dẫn nguồn cụ thể. Retrieval-Augmented Generation (RAG) được sử dụng nhằm bổ sung khả năng truy xuất tri thức từ kho tài liệu cho chatbot.

RAG giữ vai trò là phân hệ tri thức động, hỗ trợ trả lời các câu hỏi như:
\begin{itemize}
    \item Năm 2025 trường tuyển sinh những ngành nào?
    \item Ngành Công nghệ thông tin xét tuyển bằng những tổ hợp nào?
    \item Học phí dự kiến của từng ngành là bao nhiêu?
    \item Phương thức xét tuyển học bạ cần điều kiện gì?
    \item Thời gian đăng ký xét tuyển là khi nào?
\end{itemize}

Nếu không tìm thấy thông tin trong tài liệu, hệ thống cần từ chối trả lời hoặc gợi ý người dùng liên hệ bộ phận tuyển sinh. Nhờ đó chatbot không chỉ trả lời theo kịch bản cố định mà còn có khả năng khai thác tri thức cập nhật theo từng năm.
